\subsubsection{Hausdorff continuity does not provide face parameterizations}
\begin{saresult}{A.118}{Refuted}{The missing local parameterization assertion fails for polygonal disks}
There is a semialgebraic Hausdorff-continuous family of strict cone
neighborhoods of the origin in a fixed square, each a PL disk, whose
outer faces admit no jointly continuous local parameterization at one
parameter value. Both the disks and their outer faces are Hausdorff
continuous. Thus Hausdorff continuity cannot replace the joint-face
hypothesis of A.117. This does not decide whether the six-torus loop
can be filled in the weaker Hausdorff-only model.
\end{saresult}
\begin{proof}
Put $X=[-3,3]^2$, and for $0<e\leq1/4$ let
\begin{align*}
 C_e={}&([-2,2]\times[-2,2])\cup([1,2]\times[2,2+e])\\
       &\hspace{27mm}\cup([-1,2]\times[2+e,2+2e]),
 \qquad C_0=[-2,2]^2.
\end{align*}
The successive rectangles are attached along boundary intervals, so
each $C_e$ is a PL disk containing the origin in its interior. To spell
out this elementary disk assertion, a boundary interval of a polygonal
disk can be made one side of a polygonal rectangle by subdividing its
boundary and extending the boundary PL homeomorphism radially in a
triangulated disk. Attaching a second rectangle along that side gives
another disk. Apply this twice. A PL disk with a specified interior
point is PL homeomorphic to a cone on its boundary with that point as
vertex: choose any disk homeomorphism, move its interior preimage to
the center by a PL homeomorphism supported in the interior, and cone
the resulting boundary map. Hence these are strict cone neighborhoods
with the required vertex and outer face.

For $e>0$ the outer face $S_e=\partial C_e$ is the simple polygon with
successive vertices
\[
 (-2,2),\ (1,2),\ (1,2+e),\ (-1,2+e),\
 (-1,2+2e),\ (2,2+2e),\ (2,-2),\ (-2,-2).
\]
The only changes from $S_0=\partial[-2,2]^2$ lie in the strip
$2\leq y\leq2+2e$. Directly from the displayed rectangles and edges,
\[
 d_H(C_e,C_0)\leq2e,\qquad d_H(S_e,S_0)\leq2e.
\]
For positive parameters the same vertex descriptions vary continuously,
so both families are Hausdorff continuous everywhere. Their total
sets are finite unions of polynomial inequalities and therefore
semialgebraic. The exact finite certificate
\texttt{code/check\_A118\_hairpin.py} checks the polygon intersections,
the rectangular triangulation, every vertex link, and its single
boundary component at $e=1/4$. For every other positive $e$, the
piecewise-linear ambient homeomorphism that fixes $y\leq2$ and sends
$y\geq2$ to $2+4e(y-2)$ transports that certificate to $C_e$.

Suppose a joint parameterization
$E:S^1\times[0,\delta]\to X$ existed, with $E_e:S^1\to S_e$ a
homeomorphism for every $e$. Set
\[
 P_e=(0,2),\qquad Q_e=(0,2+e),\qquad P=(0,2)\in S_0.
\]
Both points tend to $P$. Compactness and injectivity of $E_0$ imply
that their preimages under $E_e$ both tend to
$a=E_0^{-1}(P)$. Indeed any convergent subsequence of either preimage
has image $P$ under $E_0$ and hence has limit $a$.
Choose a small closed parameter arc $I$ with $a$ in its interior
and $E_0(I)$ contained in the ball of radius $1/4$ about $P$.
Joint continuity makes $E_e(I)$ lie in the ball of radius $1/2$
for all sufficiently small $e$, while both preimages lie in the
interior of $I$. The subarc of $I$ joining them maps to one of the
two arcs of $S_e$ joining $P_e$ to $Q_e$. One of those arcs contains
$(1,2)$; the other contains $(-1,2+e)$. Each therefore has a point
at distance at least one from $P$. This contradicts the radius
$1/2$ bound.

Over an arbitrary real closed field the same conclusion holds against
joint definable parameterizations. Definable compactness and curve
selection replace the subsequence argument, and uniform continuity
on the compact parameter arc gives the same contradiction. No
oscillation or infinite family of combinatorial types is involved.
\end{proof}
